\documentclass{situationist-preprint}

\title{Supervaluation and the Problem of the Many}
\msid{MET-05-b}
\msfamily{MET-05}
\msdiscipline{Metaphysics}
\mstype{Companion manuscript}
\mskeywords{problem of the many; supervaluation; epistemicism; existential witnesses; maximality; higher-order vagueness}

\begin{document}
\maketitle

\begin{abstract}
MET-05-a solves the problem of the many by supervaluation: on every admissible precisification of \enquote{cloud} exactly one candidate is a cloud, so that \enquote{there is exactly one cloud} is supertrue, though no candidate is supertruly the cloud. Two objections to the supervaluationist step remain: that it asserts an existential without a witness, and that the maximality constraint on which it relies is itself vague. We answer both by the reading MET-05-a gives of supervaluation in the light of MET-18: supertruth measures what can be known, while on the epistemic account there is a fact about which candidate is the cloud. The witnessless existential is ordinary knowledge without a known witness, and the vagueness of maximality leaves the count untouched.
\end{abstract}

\section{The missing witness}

Corollary~2.3 of MET-05-a says that \enquote{there is exactly one cloud in the region} is supertrue while, for no candidate $c$, \enquote{$c$ is the cloud} is supertrue. Critics of supervaluationism, in the form given it by \citet{Fine1975}, have found this pattern intolerable when supertruth is taken as truth: an existential sentence is true although no instance of it is \citep[ch.~5]{Williamson1994}. If truth is supertruth, the semantics of the existential quantifier departs from its standard clause, and the solution of the problem of the many is bought at the price of a deviant logic of \enquote{there is}.

MET-05-a does not take supertruth to be truth. Following MET-18, it takes the use of \enquote{cloud} to fix one intended precisification, on which truth is defined, and it reads supertruth as truth on every interpretation that speakers cannot discriminate from the intended one, that is, as knowable truth. On this reading the pattern of Corollary~2.3 is not a pattern of truth values but a pattern of knowledge. It is known that there is exactly one cloud, and of no candidate is it known that it is the cloud. The existential has a witness, the candidate counted by the intended precisification; what it lacks is a known witness.

That pattern is familiar and unproblematic. We know that someone won the lottery before we know who; we know that one of the clocks in the hall is fast without knowing which. Knowledge of an existential without knowledge of any instance is the normal condition of knowledge about large or fine-grained domains. The problem of the many is a case of it, distinguished only by the source of the ignorance: we cannot know which candidate is the cloud because our use of \enquote{cloud} does not discriminate between interpretations that differ by a few droplets at the edge, and so the margin for error that knowledge requires includes them all \citep{Williamson1994}.

\section{The vagueness of maximality}

The second objection concerns the principle that does the work in MET-05-a. Maximality (P1 of MET-05-a) says that on every precisification no two distinct clouds share most of their droplets. But \enquote{most} is vague, and so, the objector says, the principle has borderline applications, and among the precisifications there may be some on which two candidates overlap in a proportion that is a borderline case of \enquote{most}, so that both are clouds and the count fails.

The objection does not touch the case. The candidates of MET-05-a (Definition~2.1) all contain every droplet that is clearly part of the cloud, and they differ only at the edge, where the droplets are sparse. Their overlap is therefore not a borderline case of \enquote{most} but a clear one: two candidates share the whole of the clearly included mass and differ by a fringe that is a small fraction of it. Every admissible sharpening of \enquote{most} counts that overlap as most, and the argument of Proposition~2.2 of MET-05-a goes through on each. Higher-order vagueness, the vagueness of \enquote{admissible precisification}, would matter only if the candidates' overlap fell near the boundary of \enquote{most}, which is the situation of two clouds of comparable size merging, not of one cloud with a ragged edge.

On the epistemic account, \enquote{most} has a sharp boundary like every vague term, and Maximality holds on the intended interpretation as it holds on all the others. What the supervaluationist model adds is that it is known to hold, since it holds on every interpretation in the margin. The reconciliation of MET-05-a thus extends to the principle as well as to its consequence: there is a fact about which candidate is the cloud, the count of one is known, and the identity of the cloud is not.

\begin{objection}[Lewis's alternative]
It would be simpler to say, with \citet{Lewis1993}, that there are many clouds, each almost identical to the others, and that we count them as one.
\end{objection}

\begin{reply}
The many candidates exist, as MET-06 requires: each is a composite whose members share the common fate of the mass. What does not follow is that each is a cloud. The supervaluationist model, read epistemically, shows that the sentence of the counting practice, \enquote{there is exactly one cloud}, can be kept true and known without counting by almost-identity, which is not counting \citep{McGee2000}. A theory that keeps the practice's sentences true is to be preferred to one that keeps them only as approximations.
\end{reply}

\begin{thebibliography}{9}\small
\bibitem[Fine(1975)]{Fine1975} Fine, K. (1975). Vagueness, truth and logic. \emph{Synthese}, 30, 265--300.
\bibitem[Lewis(1993)]{Lewis1993} Lewis, D. (1993). Many, but almost one. In J.~Bacon, K.~Campbell and L.~Reinhardt (eds.), \emph{Ontology, Causality and Mind: Essays in Honour of D.~M. Armstrong}. Cambridge University Press.
\bibitem[McGee and McLaughlin(2000)]{McGee2000} McGee, V. and McLaughlin, B. (2000). The lessons of the many. \emph{Philosophical Topics}, 28(1), 129--151.
\bibitem[Williamson(1994)]{Williamson1994} Williamson, T. (1994). \emph{Vagueness}. Routledge.
\end{thebibliography}

\end{document}
